\chapter {Rezolvare ''Cinci Pâini''}
A aduce 3 pâini, iar B aduce 2 pâini. C plătește 5 lei.  Se împarte
fiecare pâine în 3 părți egale: $3 \times 3 + 2 \times 3 = 15$
bucăți. Fiecare mesean capătă câte 5 bucăți. Reiese că A i-a dat lui C
$3 \times 3 - 5 = 4$ pâini, iar B i-a dat $2 \times 3 - 5 = 1$
pâine. Împărțeală cinstită este deci 4 lei pentru A, și 1 leu pentru B.

\chapter{Alte povești și basme}

 Vezi tabela \ref{altebasme} cu alte povești și autorii/culegătorii lor.
\lstinputlisting[language=Matlab]{Cod/Cod.m}


\begin{table}
\begin{tabular}[t]{c|p{3cm}|p{3cm}|p{3cm}|c}
\_ & Dănilă prepeleac & Prâslea cel voinic și merele de aur & Făt-Frumos din lacrimă & Greuceanu \\
\hline
 & Ion Creangă & Petre Ispirescu & Mihai Eminescu & Petre Ispirescu \\
\end{tabular}
\label{altebasme}
\caption{Alte basme.}
\end{table}